\documentclass[11pt,a4paper]{article}
\usepackage[T1]{fontenc}
\usepackage{lmodern}
\usepackage[latin1]{inputenc}
\usepackage{a4wide}
\usepackage[dvips]{graphicx}
\usepackage{float}

\usepackage[
pdfauthor={},
pdftitle={},
pdfcreator={pdftex},
]{hyperref}

\usepackage{sectsty}
\allsectionsfont{\sffamily}

\usepackage{fancyheadings} 
\pagestyle{fancy} 
\lhead{\textsf{\textbf{ACE} \\ \small{a collaborative editor}}}
\chead{}
\rhead{
\parbox[c]{3cm}{\includegraphics[height=0.875cm,width=3cm]{../../images/logo_BFH.eps}}
\parbox[c]{2.2cm}
{\tiny{\textsf{Berner Fachhochschule \\
Hochschule f�r \\
Technik und Informatik}}}}
\lfoot{}
\cfoot{\textsf{\thepage}}
\rfoot{}
\setlength{\headrulewidth}{0.6pt}
\setlength{\footrulewidth}{0.6pt}
\setlength{\topmargin}{-50pt}
\addtolength{\headheight}{50pt}

\usepackage{colortbl}

\newcommand{\headercol}[2]{\multicolumn{1}{|>{\bfseries\columncolor[gray]{0.82}}p{#1}|}{\textsf{#2}}}
\newcommand{\ace}[0]{\emph{ACE }}



\begin{document}
\setlength{\parindent}{0pt}

\begin{titlepage}
\thispagestyle{empty}
  \includegraphics[height=1.5in]{../images/pix.eps}

  \begin{center}

    {\fontsize{40}{45} \textbf{\textsf{ACE}}} \\
    \textsf{a collaborative editor} \\
        
    \vspace{36pt}
        
    {\huge{\textbf{\textsf{}}}} \\

    \vspace{36pt}

	\textsf{Berne University of Applied Sciences} \\
    \textsf{School of Engineering and Information Technology} \\
    
  \end{center}

  \vfill
  
  \begin{tabular}{ll}
   \hline

   \\

   \multicolumn{1}{>{\bfseries}p{1.5in}}{\textsf{Date:}} &
   \multicolumn{1}{>{}p{4.3in}}{\textsf{08.11.2005}}          \\
   
   \\
   
   \multicolumn{1}{>{\bfseries}p{1.5in}}{\textsf{Version:}}     &   
   \multicolumn{1}{>{}p{4.3in}}{\textsf{0.1}}                 \\

   \\
   
   \multicolumn{1}{>{\bfseries}p{1.5in}}{\textsf{Projectteam:}}                 &
   \multicolumn{1}{>{}p{4.3in}}{\textsf{Mark Bigler (biglm2@hta-bi.bfh.ch)}}  \\
   \multicolumn{1}{>{\bfseries}p{1.5in}}{}                                      &
   \multicolumn{1}{>{}p{4.3in}}{\textsf{Simon Raess (rasss@hta-bi.bfh.ch)}}    \\
   \multicolumn{1}{>{\bfseries}p{1.5in}}{}                                      &
   \multicolumn{1}{>{}p{4.3in}}{\textsf{Lukas Zbinden (zbinl@hta-bi.bfh.ch)}} \\   
   
   \\
   
   \multicolumn{1}{>{\bfseries}p{1.5in}}{\textsf{Receivers:}}                       &
   \multicolumn{1}{>{}p{4.3in}}{\textsf{Jean-Paul Dubois (doj@hta-bi.bfh.ch)}}       \\
   \multicolumn{1}{>{\bfseries}p{1.5in}}{}                                          &
   \multicolumn{1}{>{}p{4.3in}}{\textsf{Claude Fuhrer (frc@hta-bi.bfh.ch)}}       \\

   \\
   
   \multicolumn{1}{>{\bfseries}p{1.5in}}{\textsf{Location:}}               &   
   \multicolumn{1}{>{}p{4.3in}}{\textsf{Subversion Repository}} \\

   \\  
   
   \hline
  \end{tabular}

\end{titlepage}

\newpage

\tableofcontents
\newpage

\listoftables
\listoffigures
\newpage


\section{Introduction}
\ace will need a network layer that provides the following two basic features:

\begin{itemize}
 \item discovery
 \item communication
\end{itemize}

Let us explain these two basic points in greater detail.

\subsection{Discovery}
In \ace a user will be able to share a document so others can join him editing this document. Without a way of discovering other shared documents, the application is pretty useless. The discovery should be as simple as possible. There are several technologies that allow automatic discovery of services, for instance \emph{Jini}, \emph{Bonjour} (also known as zero-conf networking) and diverse peer-to-peer frameworks.

\subsection{Communication}
This is the basic requirement. Without some form of communication collaboration is infeasible. There are different technologies ranging from plain socket communication over remote method invocation to complex peer-to-peer frameworks.



\section{JGroups}
\emph{JGroups} is a toolkit for reliable multicast communication. It can be used to create groups of processes whose members can send messages to each other. The main features include:

\begin{itemize}
 \item Group creation and deletion (members can spread across LANs or WANs).
 \item Joining and leaving groups.
 \item Membership detection and notification about joined/left/crashed members
 \item Detection and removal of crashed members
 \item Sending and receiving of member-to-group messages (point-to-multipoint)
 \item Sending and receiving of member-to-member messages (point-to-point)
\end{itemize}

One of the most powerful features of \emph{JGroups} is its flexible protocol stack, which allows developers to adapt it to exactly match the application requirements and network characteristics. The benefit of this is that one has only to pay for what is actually used. By mixing and matching protocols, various differing application requirements can be satisfied. \emph{JGroups} comes with a number of protocols, for example

\begin{itemize}
 \item Transport protocols: UDP (IP multicast), TCP, JMS
 \item Fragmentation of large messages
 \item Reliable unicast and multicast message transmission. Lost messages are
       retransmitted.
 \item Failure detection: crashed members are excluded from membership
 \item Ordering protocols: atomic (all-or-none message delivery), FIFO, causal,
       total order (sequencer or token based)
 \item Membership
 \item Encryption
\end{itemize}

\emph{JGroups} allows developers to create reliable multicast applications where
\emph{reliability} is a deployment issue and does not have to be implemented by the application developer. This saves application developers significant amounts of time and allows for the application to be deployed in different environments without having to change code.

\subsection{Basics}
In order to join a group and send messages, a process has to create a channel. A channel is similar to a socket. When a client connects to a channel, it gives the name of the group it would like to join. A channel is always associated with a particular group (in it's connected state). The protocol stack takes care that channels within the same group find each other. Whenever a client connects to a channel with group name G, then it tries to find existing channels with the same name and joins them. If no members exist, a new group will be created.

When a channel is created, it is first in the unconnected state. The clients connects to the channel by calling \texttt{connect} supplying the group name that it wants to join. Once the channel is in the connected state, messages can be sent/received. A channel can be disconnected from a group by calling \texttt{disconnect} and it can be closed by calling \texttt{close}. A closed channel cannot be used anymore. Any attempt to do so results in an exception. The channel can only be connected to one group at the time. To communicate with multiple groups, multiple channels have to be created.

\subsubsection{Creating a channel}
To create a channel we can use the public constructor of \texttt{JChannel}. It looks as follows:

\begin{verbatim}
public JChannel(Object properties) throws ChannelException
\end{verbatim}

The properties arguments defines the composition of the protocol stack, that is the number and types of layers, their parameters and their order. For \texttt{JChannel} this has to be a string. For details about the composition of this string, please refer to the \emph{JGroups} users guide.

\subsubsection{Sending messages}
One of the two \texttt{send} methods of a channel are used to send messages to group members. The message payload can be any serializable object.

\begin{verbatim}
JChannel channel = ...;
channel.send(null, null, "test");
\end{verbatim}

The above code fragment sends the string "test" to all the other members in the group.

\subsubsection{Receiving messages}
The \texttt{receive} method is used to receive messages. A channel receives messages asynchronously from the network and stores them in a queue. When receive is called, the first message in the queue is removed and returned. If there is no message in the queue, the method invocation blocks.

There is a so-called \emph{building block} that provides a listener interface for receiving messages. The building block is called \texttt{PullPushAdapter}. This relieves the application developer from the task to create a separate reception thread. 

\subsection{Analysis}
\emph{JGroups} provides a very simple API to create groupware applications. A shared document in \ace could be represented as a separate group in \emph{JGroups}. 

\subsubsection{Shortcomings}
To join a group, one has to know the name of the group to join. There would have to be some mechanism to discover the available groups, or more specifically, the available groups that represent shared documents in \ace. A central server could be used to register group names or a well-known group could be used to advertise/query for available documents.

The current implementation of the group membership service is very simplistic. Joining a group is always successul. There is no way to forbid a particular member the access the group. This is a serious problem as the confidentiality of the document content cannot be guaranteed. A custom group membership service protocol implementation could probably solve this particular problem.

The \emph{JGroups} library is not programming language independent. It uses serialized objects as message payload. This can be seen as and advantage or as a disadvantage depending on the requirements of the application.


\section{Rendezvous / Bonjour}
\emph{Bonjour}, formerly known as \emph{Rendezvous}, enables automatic discovery of computers, devices, and services on IP networks. It uses industry standard IP protocols to allow devices to automatically discover each other without the need to enter IP addresses or configure DNS servers. It is a technology developed by Apple that is submitted to the IETF as part of an ongoing standards-creation process. The technology is more generally known as zero-conf networking.

There are \emph{Java} libraries available that allow to use this technology in any \emph{Java} application. The library needs a native library that is available for all major platforms. 

DNS service discovery is a way of using standard DNS programming interfaces, servers, and packet formats to browse the network for services. It is compatible with, but not dependent on, multicast DNS. Multicast DNS is a way of using familiar DNS programming interfaces, packet formats and operating semantics, in a small network where no conventional DNS server has been installed.

\subsection{Prototype}
The technology is best described by a simple prototype. We've created a simple echo server and echo client that show off the basic features of this technology.

\subsubsection{Echo Server}
The echo server is implemented in the class \texttt{ch.iserver.ace.net.bonjour.EchoServer}. It is very similar to normal echo server implementation. First a \texttt{ServerSocket} is created listening on a particular socket. Then the service is registered with DNS service discovery system represented by the \texttt{com.apple.dnssd.DNSSD} class. A listener (\texttt{com.apple.dnssd.RegisterListener}) is used to inform the application when the registration was successull. Once the registration was successul, the echo server starts to service clients as any echo server would do.

\subsubsection{Echo Client}
The echo client (see \texttt{ch.iserver.ace.net.bonjour.EchoClient}) is a simple GUI that allows to send echo requests to echo servers. It browses the network for available echo servers and lists them in a combo box. The user can select one of those echo servers and send an echo request to them.

The method \texttt{browse} of the \texttt{DNSSD} class is used to issue a browse request to the DNS service discovery system. We have to supply a \texttt{com.apple.dnssd.BrowseListener} implementation that gets called back whenever a service of specified type is found.

Whenever the user sends a message to the currently selected echo server by entering a message in the bottom textfield and hitting enter, the service has to be resolved. This happens by calling \texttt{resolve} on the \texttt{DNSSD} class. The application is called back over a \texttt{com.apple.dnssd.ResolveListener}. The method \texttt{serviceResolved} is called on this listener when the service is resolved. This provides us with the service name, host name and port of the service. As a last step we have to get the IP address of the service. This is done by calling \texttt{queryRecord} on the \texttt{DNSSD} class. We get called back by a method call to \texttt{queryAnswered} in the \texttt{com.apple.dnssd.QueryListener} we supplied to the \texttt{queryRecord} call.

At that point we do have enough information to contact the echo server and send the echo request over a plain socket connection. That's one important point to observer. This technology enables us to discover services on the network but does not impose any limitations on the way we communicate from client to server.

\section{Jini}

\section{JXTA}

\section{BEEP}
\label{sect:beep}

\emph{BEEP} stands for Blocks Extensible Exchange Protocol. It is a protocol framework for building application protocols. Many Internet protocols reinvent a set of basic functions. The most common include:

\begin{itemize}
 \item Framing: separating each request from the next
 \item Matching responses to requests
 \item Pipelining: permitting multiple outstanding requests
 \item Multiplexing: permitting multiple asynchronous requests
 \item Reporting errors
 \item Negotiating encryption
 \item Negotiating authentication
\end{itemize}

\emph{BEEP} specifies reusable tools for all these functions, instead of requiring the same decisions to be made over again for each new application. It provides a framework that integrates existing Internet standards for encryption and authentication and new standards for connection management. A networked application needs merely to supply the \emph{important} part, the part that distinguishes it from other applications, as a \emph{BEEP} profile. The mundane parts are the same for all protocols, and therefore can be coded as a library, freeing the application designer to focus on the interesting bits.

\paragraph{Why BEEP is Good:}
\begin{itemize}
 \item It solves once a number of problems common to most connection-oriented application protocols.
 \item It allows a single application to operate in multiple roles, and even multiple roles simultanously.
 \item It is efficient and requires only a small overhead.
 \item It is low-risk because it is grounded in established Internet practice and it is standardized by the IETF.
\end{itemize}

\paragraph{When BEEP is Bad:} 
\begin{itemize}
 \item If a protocol already exists and is widely deployed, there is little sense in rewriting it to take advantage of \emph{BEEP}.
 \item If a protocol is not unicast, \emph{BEEP} is an inappropriate choice of transport (for instance multicast video streaming).
 \item \emph{BEEP} is inappropriate for real-time protocols.
\end{itemize}



\subsection{Concepts}
\emph{BEEP} is a peer-to-peer protocol in the sense that there is no notion of client or server. For convenience we'll refer to the peer that starts a connection as the \emph{initiator}, and the peer accepting the connection as the \emph{listener}. When a connection is established between the two, a \emph{BEEP} session is created.

\subsubsection{Channels}
All communication in a session happens within one or more \emph{channels}. The peers require only one IP connection, which is then multiplexed to create channels. The nature of communication possible within that channel is determined by the \emph{profiles} it supports (each channel may have one or more).

The first channel, channel 0, has a special purpose. It supports the \emph{BEEP} management profile, which is used to negotiate the setup of further channels. The supported profiles determine the precise interaction between the peers in a particular channel. Defining a protocol in \emph{BEEP} comes down to the definition of profiles.

\subsubsection{Profiles}
After the establishment of a session, the initiator asks to start a channel for the particular profile or set of profiles it whishes to use. If the listener supports the profile(s), the channel will be created. Profiles themselves take one of two forms: those for initial tuning and those for data exchange.

Tuning profiles, set up at the start of communication, affect the rest of the session in some way. For instance, requesting the TLS profile ensures that channels are encrypted using Transport Layer Security. Other tuning profiles perform steps such as authentication.

Data exchange profiles set expectations between the two peers as to what sort of exchanges will be allowed in a channel, similar to the was Java interfaces set expectations between interacting objects as to what methods are available. A profile is identified by a URI.

\subsubsection{Types of Data}
\emph{BEEP} puts no limits on the kind of data a channel can carry. It uses the MIME standard to support payloads of arbitrary type.



\subsection{Prototype}
\subsubsection{Echo Server}
The \emph{BEEP} implementation of an echo server can be found in \texttt{ch.iserver.ace.net.beep.EchoServer}. The code is straighforward. First a profile registry (\texttt{org.beepcore.beep.core.ProfileRegistry}) has to be created. All the supported profiles (\texttt{org.beepcore.beep.profile.Profile}) have to be added to the created profile registry by calling \texttt{addStartChannelListener}). The profiles are registered with an URI that is used as a key to specify a specific profile when creating a new channel.

Once the registry is set up and the profiles added, the server can start listening for client connections. The class \texttt{org.beepcore.beep.transport.tcp.TCPSessionCreator}'s \texttt{listen} method is used for this purpose. The server does not need to do anything special with the returned session as all the communication is handled by the selected profile.

The echo server uses the existing \texttt{org.beepcore.beep.profile.echo.EchoProfile} profile. This profile implements the echo behavior.

\subsubsection{Echo Client}
The echo client (see \texttt{ch.iserver.ace.net.beep.EchoClient}) connects to the echo server and sends echo requests using the BEEP echo profile. First a session has to be created to the server. The method \texttt{initiate} in \texttt{org.beepcore.beep.transport.tcp.TCPSessionCreator} is used for that purpose. Once we have a session, we can start a channel by specifying the desired profile URI in a call to \texttt{startChannel} on the session.

The \texttt{sendMSG} method on the channel allows us to send a message on the channel. We need to specify a \texttt{org.beepcore.beep.core.OutputDataStream} and a \texttt{org.beepcore.beep.core.ReplyListener}. The \texttt{org.beepcore.beep.lib.Reply} implements this interface and allows to retrieve the reply from the server.



\subsection{Analysis}
\emph{BEEP} provides a framework for creating new application protocols. It fulfills our communication requirement pretty well, but not the discovery requirement. Other technologies could be used to discover shared documents, for instance \emph{Bonjour}.

\subsubsection{Positive Points}
\emph{BEEP} makes it relatively easy to create new protocols.

The messages exchanged in channels is completely up to the application. It could be serialized Java objects (whereby one would loose language independence), XML messages or some binary format. There are implementations for Java, C, TCL, Ruby and others. So \emph{BEEP} would allow us to achieve interoperability with other collaborative editors programmed in other programming languages (unless, of course, we would use serialized Java objects as messages).

\subsubsection{Negative Points}
Creating \emph{BEEP} profiles is a low-level task. One has to deal with streams, bits and bytes. 

We would need to develop a protocol specification. All the messages and possible responses would need to be defined (e.g. as XML messages).

\subsubsection{Open Points}
There are no open points, besides defining the protocol and its messages.



\subsection{APEX} 
\emph{APEX} is a \emph{BEEP} profile. On top of \emph{BEEP} it adds service discovery, application-level addressing, presence information, and permission management. In other words, it defines how to find a person with whom you are interested in communicating, regardless of where they are, not unlike an electronic mail address does. It also allows discovery of presence information, meaning that one can be notified when a coworker becomes available for a video conference. Finally, it provides a standard mechanism for permissions, specifying a service and database for defining permissions associated with user names, applications and transactions.

From the general description this sound very much like what we need. It covers both discovery and communication. Unfortunately there does not seem to be a \emph{Java} implementation for the \emph{APEX} profile.



\subsection{Resources}
\begin{itemize}
 \item \url{http://www.beepcore.org/} - official beepcore website
 \item \url{http://www-106.ibm.com/developerworks/xml/library/x-beep/index.html} - a bird's-eye view on beep
 \item \url{http://www-128.ibm.com/developerworks/xml/library/x-beep2.html} - a worm's eye view on beep
\end{itemize}



\end{document}
